% ============================================================
%
%   WORKING NOTES ON CHIRAL KOSZUL DUALITY
%
%   Raeez Lorgat, 2024--2026
%
%   Compile standalone: pdflatex working_notes.tex
%
% ============================================================

\documentclass[11pt]{article}

\usepackage[T1]{fontenc}
\usepackage[utf8]{inputenc}
\usepackage{lmodern}
\frenchspacing

\usepackage[cmintegrals,cmbraces,noamssymbols]{newtxmath}
\usepackage{ebgaramond}

\usepackage[
  activate={true,nocompatibility}, final,
  tracking=true, kerning=true, spacing=true,
  factor=1100, stretch=10, shrink=10
]{microtype}

\usepackage{mleftright}\mleftright

\let\Bbbk\undefined \let\openbox\undefined
\usepackage{amsmath,amssymb,amsthm}
\usepackage{mathtools}
\usepackage{enumitem}
\usepackage[margin=1.15in]{geometry}
\usepackage[unicode]{hyperref}
\usepackage{tikz-cd}

\setlength{\emergencystretch}{3em}
\raggedbottom

% --- Environments ---
\usepackage{thmtools}
\declaretheoremstyle[
  spaceabove=\topsep, spacebelow=\topsep,
  headfont=\normalfont\scshape,
  notefont=\normalfont\itshape,
  bodyfont=\normalfont,
  postheadspace=0.5em, headpunct={.}
]{notesthm}
\declaretheoremstyle[
  spaceabove=\topsep, spacebelow=\topsep,
  headfont=\normalfont\itshape,
  notefont=\normalfont\itshape,
  bodyfont=\normalfont,
  postheadspace=0.5em, headpunct={.}
]{notesdef}

\declaretheorem[style=notesthm, name=Theorem, numberwithin=section]{theorem}
\declaretheorem[style=notesthm, name=Lemma, sibling=theorem]{lemma}
\declaretheorem[style=notesthm, name=Proposition, sibling=theorem]{proposition}
\declaretheorem[style=notesthm, name=Corollary, sibling=theorem]{corollary}
\declaretheorem[style=notesdef, name=Definition, sibling=theorem]{definition}
\declaretheorem[style=notesdef, name=Remark, sibling=theorem]{remark}
\declaretheorem[style=notesdef, name=Observation, sibling=theorem]{observation}
\declaretheorem[style=notesdef, name=Warning, sibling=theorem]{warning}
\declaretheorem[style=notesdef, name=Example, sibling=theorem]{example}
\declaretheorem[style=notesdef, name=Conjecture, sibling=theorem]{conjecture}
\declaretheorem[style=notesdef, name=Question, sibling=theorem]{question}

% --- Macros ---
\newcommand{\cA}{\mathcal{A}}
\newcommand{\cB}{\mathcal{B}}
\newcommand{\cC}{\mathcal{C}}
\newcommand{\cM}{\mathcal{M}}
\newcommand{\cN}{\mathcal{N}}
\newcommand{\cO}{\mathcal{O}}
\newcommand{\cW}{\mathcal{W}}
\newcommand{\cH}{\mathcal{H}}
\newcommand{\cZ}{\mathcal{Z}}
\newcommand{\cD}{\mathcal{D}}
\newcommand{\cI}{\mathcal{I}}
\newcommand{\fg}{\mathfrak{g}}
\newcommand{\fh}{\mathfrak{h}}
\newcommand{\fn}{\mathfrak{n}}
\newcommand{\fz}{\mathfrak{z}}
\newcommand{\fsl}{\mathfrak{sl}}
\newcommand{\barB}{\overline{B}}
\newcommand{\B}{\bar{B}}
\newcommand{\C}{\mathbb{C}}
\newcommand{\Z}{\mathbb{Z}}
\newcommand{\Q}{\mathbb{Q}}
\newcommand{\bP}{\mathbb{P}}
\newcommand{\bE}{\mathbb{E}}

\newcommand{\Ass}{\mathsf{Ass}}
\newcommand{\Com}{\mathsf{Com}}
\newcommand{\Lie}{\mathsf{Lie}}
\newcommand{\Eone}{\mathsf{E}_1}
\newcommand{\Etwo}{\mathsf{E}_2}
\newcommand{\Ethree}{\mathsf{E}_3}
\newcommand{\Einf}{\mathsf{E}_{\infty}}
\newcommand{\Linf}{\mathsf{L}_{\infty}}
\newcommand{\Ainf}{\mathsf{A}_{\infty}}
\newcommand{\Pinf}{\mathsf{P}_{\infty}}
\newcommand{\chirAss}{\mathsf{Ass}^{\mathrm{ch}}}
\newcommand{\chirCom}{\mathsf{Com}^{\mathrm{ch}}}
\newcommand{\chirLie}{\mathsf{Lie}^{\mathrm{ch}}}
\newcommand{\Ran}{\operatorname{Ran}}
\newcommand{\Res}{\operatorname{Res}}
\newcommand{\FM}{\mathrm{FM}}
\newcommand{\CE}{\mathrm{CE}}
\newcommand{\DS}{\mathrm{DS}}
\newcommand{\MC}{\mathrm{MC}}
\newcommand{\ad}{\operatorname{ad}}
\newcommand{\id}{\mathrm{id}}
\newcommand{\Op}{\operatorname{Op}}
\newcommand{\Fun}{\operatorname{Fun}}
\newcommand{\Hom}{\operatorname{Hom}}
\newcommand{\Ext}{\operatorname{Ext}}
\newcommand{\Loc}{\operatorname{Loc}}
\newcommand{\Defcyc}{\operatorname{Def}_{\mathrm{cyc}}}
\newcommand{\dzero}{d_0}
\newcommand{\dfib}{d_{\mathrm{fib}}}
\newcommand{\Dg}[1]{D_{#1}}
\newcommand{\Dtot}{D_{\mathrm{tot}}}
\newcommand{\normord}[1]{:\!#1\!:}
\newcommand{\Fact}{\mathsf{Fact}}
\newcommand{\thick}{\operatorname{thick}}

\DeclareMathOperator{\gr}{gr}
\DeclareMathOperator{\Bl}{Bl}
\DeclareMathOperator{\Spec}{Spec}
\DeclareMathOperator{\Jac}{Jac}

% Separator
\newcommand{\sep}{\bigskip\begin{center}$*$\end{center}\medskip}

\title{\textit{Working Notes on Chiral Koszul Duality}}
\author{Raeez Lorgat}
\date{2024--2026}

\begin{document}
\maketitle

\begin{quote}
\itshape
These notes are about what happens when you compute with a
$1$-form, what it teaches you when you fail, and why the failures
contain more information than the successes.
\end{quote}

\tableofcontents

% =================================================================
\section{A 1-form}
\label{sec:one-form}
% =================================================================

\begin{equation}\label{eq:seed}
\eta_{ij} \;=\; d\log(z_i - z_j).
\end{equation}

A $1$-form on the configuration space of two distinct points on a
curve.  It has a simple pole along $z_i = z_j$, with residue $1$.
Not residue $k$, not residue depending on anything.  Just $1$.

A double pole $dz/(z-w)^2$ has no coordinate-independent residue.
A regular form extracts nothing.  Only the logarithmic derivative
threads the needle: a pole strong enough to detect, mild enough
to be canonical.

For three points:
\begin{equation}\label{eq:arnold}
\eta_{12} \wedge \eta_{23}
\;+\; \eta_{23} \wedge \eta_{31}
\;+\; \eta_{31} \wedge \eta_{12}
\;=\; 0.
\end{equation}
This is Arnold's relation.  It holds because $(z_1 - z_2)(z_2 - z_3)
+ (z_2 - z_3)(z_3 - z_1) + (z_3 - z_1)(z_1 - z_2) = 0$, which
holds because addition is commutative.  The deepest identity in the
subject follows from the fact that $a + b + c = 0$ implies
$ab + bc + ca = 0$ when $a + b + c = 0$.
It is the only relation.

\sep

Now place elements $\phi_1, \ldots, \phi_n$ of a chiral algebra
$\cA$ at $n$ points on a smooth curve $X$.  Dress them in
logarithmic forms.  Define a differential by extracting residues
at collisions.  Arnold's relation kills every triple residue.
So $d^2 = 0$, and you have a cochain complex: the \textbf{bar
complex} $\B_X(\cA)$.

Under a coordinate change $z \mapsto w(z)$, the propagator picks
up a correction:
$\eta_{ij} \mapsto \eta_{ij} +
d\log\frac{w(z_i) - w(z_j)}{z_i - z_j}$.
Near the diagonal this becomes $d\log w'(z_i)$, which is the
standard cocycle for $\mathrm{Diff}(X)$.  The propagator
transforms as a BRST ghost for diffeomorphisms.


% =================================================================
\section{The trouble with two poles}
\label{sec:trouble}
% =================================================================

The Heisenberg OPE $\alpha(z)\alpha(w) = k/(z-w)^2$ has one pole.
The bar differential extracts a scalar and stops.  Everything works.
Cohomology counts partitions.  This should be a warning.

Kac--Moody has two poles:
\[
J^a(z) J^b(w) \;=\; \frac{k\kappa^{ab}}{(z-w)^2}
  \;+\; \frac{f^{ab}{}_c\, J^c(w)}{z-w} \;+\; \mathrm{reg}.
\]
The simple pole is the Lie bracket.  The double pole is the level.
So the bar differential splits: $d = d_{\mathrm{bracket}} +
d_{\mathrm{level}}$.

$d_{\mathrm{bracket}}^2 \neq 0$ whenever
the OPE has both poles.  This holds across all $2048$ sign
conventions arising from eleven independent binary choices in the
bar construction (grading, suspension, contraction, Koszul signs,
etc.).  No convention saves it.

What happens: iterated residues at a triple collision produce cross-terms
between the bracket and the level.  Jacobi cancels the
bracket-bracket part.  The bracket-level cross-terms are
proportional to $k$ and survive.  The relation that actually holds
is
\begin{equation}\label{eq:cross-term}
d_{\mathrm{bracket}}^2
\;=\; -\{d_{\mathrm{bracket}},\, d_{\mathrm{level}}\}.
\end{equation}
The pair $(d_{\mathrm{bracket}}, d_{\mathrm{level}})$ is not a
bicomplex.  The total $d$ squares to zero by the Borcherds identity,
which is not the Jacobi identity with extra terms.  It is the
algebraic shadow of Stokes' theorem on the Fulton--MacPherson
compactification $\overline{C}_n(X)$.  At the corner where three
points collide, the two adjacent boundary faces have opposite
orientations.  The cancellation is topological.  No algebraic sign
convention produces it.

Decomposing the differential by pole order clarifies the structure
but destroys the computation.  The Borcherds identity governs all
poles at once.  The Jacobi identity governs one.  When there are
two poles, these are different identities.

\sep

Equation~\eqref{eq:cross-term} is the conformal anomaly.  The
classical BRST operator does not square to zero.  The quantum
correction restores it.  The BV quantum master equation
$\hbar\Delta S + \tfrac{1}{2}\{S,S\} = 0$ is the same equation.

The obstruction to $d^2 = 0$ is not an obstruction to the theory.


% =================================================================
\section{Three letters $d$}
\label{sec:three-d}
% =================================================================

On $\bP^1$ the propagator is single-valued and Arnold's relation
holds.  On an elliptic curve $E_\tau$ the function $z_1 - z_2$ is
not single-valued.

Replace it by the odd theta function, with a correction forced by
double-periodicity:
\[
\eta^{(1)}_{12} = d\log\vartheta_1(z_1 - z_2 \mid \tau)
  + \frac{2\pi i\,\operatorname{Im}(z_1 - z_2)}
    {\operatorname{Im}\tau}\, d\bar{z}_1.
\]
Arnold's relation fails.  Of course it does---the propagator
is no longer single-valued.  What remains is the Arakelov $(1,1)$-form:
\begin{equation}\label{eq:curvature}
\dfib^{\,2} = \kappa(\cA) \cdot \omega_g.
\end{equation}

One number, $\kappa(\cA)$, extracted from the leading OPE
coefficient, controls everything.

\sep

Three differentials act on the genus-$g$ bar complex.

\medskip
\textbf{(i)} $\dzero$: genus-$0$ collision differential.  Built
from $d\log(z_i - z_j)$ on $\bP^1$.  Squares to zero by Arnold.
This is classical Koszul duality.

\medskip
\textbf{(ii)} $\dfib$: fiberwise differential on a fixed
$\Sigma_g$.  $\dfib^{\,2} = \kappa(\cA) \cdot \omega_g \neq 0$.
The world is curved.

\medskip
\textbf{(iii)} $\Dg{g}$: total corrected differential, with
$A$-period parameters $t_k = \oint_{A_k}\omega_k$.
$\Dg{g}^2 = 0$.

The correction that restores nilpotence is a Lagrangian splitting.
The $A$-cycles form a Lagrangian subspace of $H^1(\Sigma_g, \C)$
for the intersection pairing; the curvature lies in the symplectic
complement; the $A$-period corrections project it to zero.
Nilpotence comes from the symplectic topology of the curve.

In the curved $\Ainf$ framework: $m_0 = \kappa \cdot \omega_g$,
and $m_1^2(a) = [m_0, a]$.  The condition $\kappa = 0$ is
simultaneously: $\dfib^2 = 0$; anomaly cancellation for the
bosonic string; $c = 26$; the bar complex being a genuine derived
object.  When $\kappa \neq 0$ the bar complex is curved, and the
curvature is the invariant.

At genus~$1$ the tadpole amplitude is $\kappa \cdot E_2(\tau)$.
The Eisenstein series transforms as
$E_2(-1/\tau) = \tau^2 E_2(\tau) + 6\tau/(\pi i)$, and the
anomalous term is controlled by $\kappa$.

Across all genera, the generating function of the obstruction
classes is
\[
\sum_{g \geq 1} \kappa(\cA) \cdot \lambda_g \cdot \hbar^{2g-2}
\;=\;
\kappa(\cA) \cdot \Bigl(\hat{A}(i\hbar) - 1\Bigr).
\]
The $\hat{A}$-genus, from index theory.  It has no business
being here.


% =================================================================
\section{Fourier}
\label{sec:fourier}
% =================================================================

The bar construction is a Fourier transform.

\begin{center}
\small
\begin{tabular}{lll}
\textbf{Classical Fourier} & \textbf{Chiral bar} \\
\hline
Kernel: $e^{-2\pi i x\xi}$
  & Kernel: $\eta_{ij} = d\log(z_i - z_j)$ \\[3pt]
$\widehat{f * g} = \hat{f} \cdot \hat{g}$
  & $\mathbb{D}_{\Ran} \circ \B \simeq \B \circ (-)^!$
  \quad(Thm~A) \\[3pt]
$\hat{\hat{f}}(x) = f(-x)$
  & $\Omega(\B(\cA)) \xrightarrow{\sim} \cA$
  \quad(Thm~B) \\[3pt]
$\|f\|^2 = \|\hat{f}\|^2$
  & $Q_g(\cA) \oplus Q_g(\cA^!) = H^*(\overline{\cM}_g, \cZ_\cA)$
  \quad(Thm~C) \\[3pt]
Char.\ fn.\ $\to$ moments
  & $\kappa(\cA) \to$ genus tower via $\hat{A}$
  \quad(Thm~D)
\end{tabular}
\end{center}

When $X$ is a point, the propagator vanishes, $d_{\mathrm{res}}$
reduces to the binary bracket, and you recover the classical bar
construction of Loday--Vallette.  Operadic Koszul duality is the
zero-dimensional shadow.

For the Heisenberg on $E_\tau$, the bar complex is the
Fourier--Mukai kernel on $\Jac(E_\tau) \times \Jac(E_\tau)^\vee$.
The chiral bar construction \emph{is} the Fourier--Mukai transform,
generalized from abelian to non-abelian.

The logarithm converts multiplication to addition.  The bar
construction converts the OPE (multiplicative, with poles) to the
bar differential (nilpotent, additive).  The cobar reverses this:
it inserts delta-function currents at boundary strata.  Verdier
duality on $\overline{C}_n(X)$ exchanges residues with inclusions,
so bar and cobar are Verdier dual.  Theorems A--D are four faces of
this one duality.

At $n = 4$, the compactification $\overline{M}_{0,5}$ is the
del Pezzo surface $\Bl_4(\bP^2)$, with $10$ boundary divisors.
Five respect cyclic ordering.  The relation $\partial^2 = 0$ on
this surface yields the $\Ainf$ pentagon identity.  The coherence
relations of homotopical algebra come from blowing up
$\bP^2$ four times.  Why should they?

For Koszul algebras: $m_k = 0$ for $k \geq 3$.  For interacting
algebras: $m_3 \neq 0$, and the pentagon is the first obstruction
to formality.  Whether a chiral algebra is Koszul is a question
about the geometry of a del Pezzo surface.


% =================================================================
\section{Nilpotence is the logarithm of periodicity}
\label{sec:nilpotence-log}
% =================================================================

The classical logarithm converts multiplication to addition.  The
bar construction does the same thing: it converts the OPE (a product,
with poles) to the bar differential (a sum, nilpotent).  The
last section presented this as analogy.  It is not analogy.

\sep

Write $\log^{\mathrm{cat}}(\cA) := \B_X(\cA)$ and
$\exp^{\mathrm{cat}}(\cC) := \Omega_X(\cC)$.  Then:

\medskip
\begin{tabular}{@{}ll@{}}
$\exp(\log(z)) = z$ on $|z-1| < 1$
  & $\Omega(\B(\cA)) \xrightarrow{\sim} \cA$ on the Koszul locus \\[3pt]
$\log(1{+}z) = z - z^2/2 + \cdots$
  & PBW spectral sequence: Taylor expansion of the log \\[3pt]
Convergence radius
  & $E_2$ degeneration (Koszulness) \\[3pt]
Analytic continuation past radius
  & Coderived category \\[3pt]
$\log(ab) = \log(a) + \log(b)$
  & $Q_g(\cA) + Q_g(\cA^!) = H^*(\overline{\cM}_g, \cZ_\cA)$
  \quad(Thm~C) \\[3pt]
$\log(1) = 0$
  & $\B(k)$ is acyclic \\[3pt]
Single-valued on simply connected domain
  & $\dzero^2 = 0$ at genus $0$ \\[3pt]
Multi-valued on $\C^\times$; monodromy $2\pi i$
  & $\dfib^{\,2} = \kappa \cdot \omega_g$ at genus $g \geq 1$ \\[3pt]
Residue at branch point
  & Curvature $\kappa(\cA)$ \\[3pt]
Leading coefficient of $\log(1{+}z) = z + \cdots$
  & Modular characteristic $\kappa(\cA)$ \quad(Thm~D)
\end{tabular}

\medskip

Every row is a theorem.  Not by design.

\sep

The PBW spectral sequence is the Taylor expansion.  The filtration
degree is the order of the expansion: the $p$-th filtered piece
retains contributions from at most $p$ points on $X$, the first
$p$ terms of the log series.  The differential $d_1$ extracts the
leading (binary OPE) term.  Higher differentials $d_r$ are
higher-order corrections from $r$-fold collisions.

$E_2$ degeneration means the Taylor expansion terminates: the
log is polynomial in the filtration variable.  Off the Koszul locus,
the series diverges.  The coderived category is the analytic
continuation---the enlargement of the derived category that makes the
divergent log well-defined.  Positselski's acyclicity theorem is
the statement that the analytic continuation exists.

The Koszul locus is the convergence radius.  For KM: linear OPE
poles, finite radius, convergent.  For admissible-level affine
algebras: the OPE singularity structure forces non-terminating
higher differentials, the radius shrinks to zero, and only the
completed category survives.

\sep

The genus transition is the passage from single-valued to multi-valued.

On $\bP^1$, $\log(z_1 - z_2)$ is single-valued.  Arnold's relation
holds.  $\dzero^2 = 0$.  No monodromy.  This is Theorems A and B.

On $\Sigma_g$ with $g \geq 1$, the function $z_1 - z_2$ is not
single-valued.  Its replacement $E(z_1, z_2)$, the prime form,
has monodromy $\exp(2\pi i \oint_{B_k} \omega_j)$ around the
$B$-cycles.  The categorical logarithm inherits this monodromy.
Arnold's relation acquires a correction.  $\dfib^{\,2} =
\kappa \cdot \omega_g \neq 0$.  The curvature $\kappa$ is the
residue at the branch point.

The Gauss--Manin correction restores nilpotence: $\Dg{g}^2 = 0$.
The total differential trivializes the multi-valued log on the
universal cover of moduli space.  The three differentials $\dzero$,
$\dfib$, $\Dg{g}$ are the germ, the multi-valued lift, and the
trivialization.  Three levels of the same logarithm.

\sep

Complementarity is the additivity of the logarithm.

$Q_g(\cA) + Q_g(\cA^!) = H^*(\overline{\cM}_g, \cZ_\cA)$.  The
left side is $\log(\cA) + \log(\cA^!)$.  The right side is
$\log(\cA \otimes \cA^!)$.  Theorem~C says the categorical
logarithm is a homomorphism.  The Lagrangian decomposition is the
decomposition of the total log into two branches---one for each sheet
of the multi-valued function.

The functional equation of the Riemann zeta relates $\zeta(s)$ and
$\zeta(1-s)$ through the gamma factor.  Theorem~C relates
$\log^{\mathrm{cat}}(\cA)$ and $\log^{\mathrm{cat}}(\cA^!)$ through
the center $\cZ_\cA$.  These are the same pattern: a duality that
decomposes total data into complementary halves.

\sep

The characteristic hierarchy is the order expansion of the
logarithm.

$\kappa(\cA)$: the leading coefficient.  Determines the scalar
package.  Its exponential is the period $N = 24q/\gcd(p,24)$ for
$c = p/q$.  This is $\exp(2\pi i N c/24) = 1$ --- the monodromy
of the T-matrix is the exponential of the central charge.

$\Delta_\cA$: the eigenvalue spectrum.  Determines the periodicity
profile $\Pi(\cA)$.  Each T-matrix eigenvalue $\lambda_i =
\exp(2\pi i(h_i - c/24))$ is the exponential of a conformal weight.

$\Theta_\cA$: the complete logarithm.  Its exponential is the
modular functor: the action of $\Gamma_g$ on conformal blocks.

One is the leading term of the next.  $\kappa$ is the trace of
$\Delta$, $\Delta$ is the eigenvalue decomposition of $\Theta$.
In complex analysis: the winding number, the set of periods, the
full multi-valued function.

\sep

Quantization is exponentiation.  Dequantization is taking the
logarithm.

The passage from Lie algebra to Lie group is $\exp$: nilpotent
structure (Jacobi identity) produces periodic structure
($\exp(2\pi i X) = 1$ for lattice elements).  The bar construction
reverses this: it extracts the nilpotent infinitesimal data from
the periodic algebraic structure.  The cobar exponentiates it back.

The screening charges $Q_i$ of lattice vertex algebras have
$Q_i^2 = 0$.  This is the nilpotent datum.  The quantum group
relation $e_i^N = 0$ at $q = e^{2\pi i/N}$ is the exponential
datum: the $N$-th power of the screening operator vanishes in
the quantum group.  The root of unity is the monodromy of the
categorical logarithm around the $N$-th root.

The passage $Q^2 = 0 \Rightarrow e^N = 0$ is not metaphorical
exponentiation.  The braiding matrix $R$ of $\cO_q$ satisfies
$R = \exp(r)$ where $r$ is the classical $r$-matrix---which is
the binary OPE residue in the bar differential.  The quantum group
is the Lie group of the chiral Lie algebra, and the bar construction
is the map from group to algebra.

\sep

The four theorems are the four properties of a logarithm:
\begin{itemize}[nosep]
\item A: the logarithm exists;
\item B: $\exp \circ \log = \mathrm{id}$ on the convergence domain;
\item C: the branch structure is controlled (complementary sheets);
\item D: the leading coefficient determines the local behavior.
\end{itemize}

The five master conjectures are five aspects of one question: what
is the categorical logarithm for factorization algebras on algebraic
curves?

MC1: does it converge?  MC3: what is the convergence domain?
MC4: does it extend to infinite rank?  MC5: do the physical and
mathematical logarithms coincide?  Periodicity: what is the
monodromy group?

Not five problems.  One object, five windows.


% =================================================================
\section{Two things called $\Eone$}
\label{sec:two-e1}
% =================================================================

Which $\Eone$?

On a smooth curve $X$:
\[
\Einf\text{-chiral}
\;\subset\; \Pinf\text{-chiral}
\;\subset\; \Eone\text{-chiral}.
\]

$\Einf$-chiral = vertex algebra.  Skew-symmetric bracket,
unordered configuration spaces.  Heisenberg, KM, Virasoro,
$\cW$-algebras.

$\Eone$-chiral = drop skew-symmetry.  $R$-twisted locality:
$J_a(z) J_b(w) - R^{ab}_{cd}(z{-}w)\, J_d(w) J_c(z) = 0$.
When $R = P$ (permutation): $\Einf$.  When $R \neq P$: ordered
configuration spaces.  Yangians, quantum lattice algebras.

Meanwhile, separately, in a different part of mathematics: the
little $n$-disks operad $\mathsf{E}_n$ lives on real
$n$-manifolds.  A complex curve is real $2$-dimensional.  The
chiral bar complex lives at $n = 2$ (Arnold relations).
Chern--Simons on a $3$-manifold lives at $n = 3$ (Totaro
relations).

An $\Eone$-\emph{chiral} algebra lives on a curve.  A topological
$\Eone$-algebra lives on an interval.  These are different objects
that share an overloaded name.  On an interval there are no
configuration spaces, no Arnold relation, no propagator.  The
subject does not become simpler.  It ceases to exist.

\sep

The Kazhdan--Lusztig equivalence sends $\widehat{\fg}_k$-modules
to $U_q(\fg)$-modules, with $q = e^{\pi i/(k+h^\vee)}$ and
braiding by the universal $R$-matrix.  Verdier duality on
$\overline{C}_n(X)$ reverses the braiding: $R \mapsto R^{-1}$.

Without ordering: Verdier duality acts as
$k \mapsto -k - 2h^\vee$ (the Feigin--Frenkel involution).  With
ordering: $R \mapsto R^{-1}$ (the Drinfeld--Kohno theorem).

So:
\[
Y_\hbar(\fsl_N)^{!,\mathrm{ch}}
\;\cong\;
Y_{R^{-1},\hbar}(\fsl_N)^{\mathrm{ch}}.
\]

This single identity unifies three theorems.  Feigin--Frenkel is
its $\Einf$ shadow.  Drinfeld--Kohno is its monodromy.
$\Ass^! = \Ass$ is its shadow at a point.

Lattice vertex algebras make the $\Einf$/$\Eone$ transition
visible.  A $2$-cocycle $\varepsilon$ on a lattice $\Lambda$
produces $V_\Lambda^\varepsilon$.  When $\varepsilon$ is symmetric:
$\Einf$.  Otherwise: strictly $\Eone$.  Koszul duality inverts the
cocycle: $\varepsilon \mapsto \varepsilon^{-1}$.  For even unimodular
$\Lambda$: self-duality, $V_\Lambda^! \cong V_\Lambda$.  The $E_8$
lattice vertex algebra is Koszul self-dual.

A Coisson algebra (Beilinson--Drinfeld) is a Poisson vertex algebra,
which is a commutative $\cD_X$-algebra with $\lambda$-bracket.  It
is not a chiral algebra.  The quantization tower runs:
\[
\text{Coisson}
\;\xrightarrow{\text{star product}}\;
\Einf\text{-chiral}
\;\xrightarrow{R\text{-matrix}}\;
\Eone\text{-chiral}.
\]
Holomorphic-topological gauge theories in $3$d produce Coisson
algebras.  Their quantization lands at $\Einf$, not $\Eone$.
Reaching $\Eone$ requires the KZ equations on a curve---a
$2$-dimensional phenomenon.  Three dimensions is not enough.

From $3$d: Chern--Simons on $\Sigma_g \times [0,1]$ produces the
WZW model on $\Sigma_g$.  The perturbative expansion gives $\Ainf$
structure on flat connections---the $3$d version of the curved
$\Ainf$ on the genus-$g$ bar complex.  The curvature $m_0$ in $2$d
is the CS level in $3$d.

From $4$d: the Costello--Li holomorphic-topological twist of
$\cN = 2$ gauge theory on $\C \times \C^*$, reduced along $|w|=1$,
yields a factorization algebra on $\C$.  When conformal symmetry
survives: a chiral algebra.  Otherwise: a factorization algebra
that lives in the Costello--Gwilliam framework, outside the
$\cD$-module world.


% =================================================================
\section{The BV complex is the bar complex}
\label{sec:bv}
% =================================================================

The BV bracket $\{-,-\}$ of degree $+1$ is the simple-pole OPE
residue.  The BV Laplacian $\Delta$ is the double-pole differential.
The quantum master equation $\hbar\Delta e^{S/\hbar} = 0$ is
bar-cobar Verdier compatibility.  These are not analogies.  They are identities.

The conformal anomaly is one object seen from three windows:
\begin{itemize}[nosep]
\item the curvature element $m_0 = \kappa/2$ of the curved $\Ainf$;
\item the obstruction class in $R^*(\overline{\cM}_g)$;
\item the failure $\hbar\Delta S + \tfrac{1}{2}\{S,S\} \neq 0$
  when $c \neq 26$.
\end{itemize}

For the bosonic string: $\kappa(\cH^{26}) = 13$,
$\kappa(bc) = -13$, total $\kappa = 0$.  No curvature.  No anomaly.
The QME holds.  These are the same sentence, said three ways.

The propagator transforms under coordinate changes as a BRST ghost
for diffeomorphisms.  The bar complex is a coordinate-free BRST
complex.  At genus~$0$ this is proved.  At higher genus it requires
Costello's renormalization machinery, and waits downstream.


% =================================================================
\section{Koszulness is a theorem, not a definition}
\label{sec:koszulness}
% =================================================================

$U(\fg)$ is Koszul.  This says nothing about $\widehat{\fg}_k$.
The classical bar complex uses one product.  The chiral bar complex
uses the full OPE, all poles at once.  Higher poles are not
perturbations.  The Sugawara tensor
$T = \frac{1}{2(k+h^\vee)}\sum_a \normord{J^a J^a}$ produces
quartic poles in $T(z)T(w)$ via double Wick contractions.  These
have no classical analogue.

There is a circularity.  To compute bar cohomology: assume Koszulness
(spectral sequence collapses at $E_2$).  Bar cohomology would then
verify Koszulness (via $H_\cA(t) \cdot H_{\cA^!}(-t) = 1$).
Breaking this circle requires an independent collapse proof, and
every such proof is specific to its family.

\textsc{Kac--Moody.}\enspace
Whitehead's theorem ($H^n(\fg, V) = 0$ for $n = 1,2$) handles
the off-diagonal, and Killing contraction handles the invariant
piece.  The PBW enrichment at $E_1$ is $\fg \otimes \fg$
(dimension $9$ for $\fsl_2$), not $\Lambda^2(\fg)$ (dimension $3$).

\textsc{Virasoro.}\enspace
The subalgebra $\{L_{-1}, L_0, L_1\} \cong \fsl_2$ acts
independently of the central charge, because $m^3 - m = 0$ for
$m \in \{-1, 0, 1\}$.  The Casimir eigenvalues at every weight
$h \geq 2$ are positive.

\textsc{Principal $\cW_N$.}\enspace
Unique weight-$2$ mechanism from the stress tensor; triangular
$d_2$ argument.

There is no universal collapse theorem.  Each algebra must be
convicted on its own evidence.

\sep

Free fields (Heisenberg, free fermion) have one OPE pole:
$d_{\mathrm{bracket}} = 0$, no Borcherds needed, no PBW needed.
Their bar cohomology grows sub-exponentially---partitions,
$e^{C\sqrt{n}}$.  Interacting fields (KM, Virasoro, $\cW$-algebras)
have two poles, and their bar cohomology grows
exponentially---$3^n$, $8^n$.  The jump from one pole to two is
visible in the asymptotics.  One pole: garden.  Two poles:
wilderness.


% =================================================================
\section{What the bar complex counts}
\label{sec:bestiary}
% =================================================================

\begin{center}
\small
\renewcommand{\arraystretch}{1.15}
\begin{tabular}{lcll}
\textbf{Algebra} & \textbf{Dims (1--6)}
  & \textbf{Growth} & \textbf{Sequence}\\
\hline
Heisenberg & $1,1,1,2,3,5$ & $e^{C\sqrt{n}}$
  & partitions $p(n)$ \\
Free fermion & $1,1,2,3,5,7$ & $e^{C\sqrt{n}}$
  & partitions $p(n{-}1)$ \\
$bc$ ghosts & $2,3,6,13,28,59$ & $2^n$
  & --- \\
$\beta\gamma$ & $2,4,10,26,70,192$ & $3^n/\sqrt{n}$
  & A025565 \\
$\widehat{\fsl}_2$ & $3,6,15,36,91,232$ & $3^n/n^{3/2}$
  & Riordan numbers \\
Virasoro & $1,2,5,12,30,76$ & $3^n/n^{3/2}$
  & Motzkin differences \\
$\widehat{\fsl}_3$ & $8,36,204,\mathbf{1352}$
  & $8^n$ & (conjectured) \\
$\cW_3$ & $2,5,16,\mathbf{52},\mathbf{171}$
  & $3.3^n$ & (conjectured) \\
$Y(\fsl_2)$ & $4,10,28,\mathbf{82}$
  & $3^n$ & $3^n + 1$ (conjectured)
\end{tabular}
\end{center}

The $\widehat{\fsl}_2$ generating function is
\[
H(x) = \frac{1 + x - \sqrt{(1-3x)(1+x)}}{2x(1+x)}.
\]
These are the Riordan numbers.  Nobody asked for them.
They count non-crossing partitions
with a certain forbidden pattern.  After Drinfeld--Sokolov reduction
to the Virasoro algebra, the generating function rearranges into
Motzkin differences $M(n{+}1) - M(n)$, lattice paths in the
half-plane.  Both come from the same discriminant
$\Delta = (1-3x)(1+x)$; the DS reduction is a change of variables
that preserves the branch cut.

The discriminant is a DS-reduction invariant.  The $\fsl_3$ family
has discriminant $1 - 3x - x^2$, with roots involving $\sqrt{13}$
(and $13 = 3^2 + 4$).  The recurrence depth of the generating
function equals $\operatorname{rank}(\fg) + 1$ in all known cases.

The conjectured Yangian formula $H^n(Y(\fsl_2)) = 3^n + 1$ would
give the only interacting algebra with a rational generating
function---no branch points at all.  This may reflect the fact that
Yangian OPE data is rational (Yang $R$-matrix), while KM data
involves transcendental functions (theta, Weierstrass) at higher
genus.  Seven independent consistency checks support the conjecture
through degree~$10$.  If true, it demands an explanation.

\sep

$\cW$-algebra Koszul duality for the principal nilpotent:
$\cW^k(\fg)^! \cong \cW^{k'}(\fg)$, where $k' = -k - 2h^\vee$.
The proof goes through the Wakimoto free-field resolution:
compute the bar complex of the free fields, observe that the BRST
screening momenta reverse sign ($\alpha \to -\alpha$), forcing the
dual level.

For $\fsl_2$: $k' = -k-4$, and $c(k) + c(k') = 26$.  For
$\fsl_3$: $c + c' = 100$.  The anomaly ratio
$\rho(\fg) = \kappa/c = \sum (m_i+1)^{-1}$, summed over exponents,
is level-independent.  For $\fsl_2$: $1/2$.  For $\fsl_N$:
the harmonic number $H_N - 1$.  The genus-$1$ free energy is
$\kappa/24$ for every principal $\cW$-algebra.

For non-principal nilpotent orbits, the expected duality is
$\cW^k(\fg,f)^! \cong \cW^{k'}(\fg, f^D)$, where $f^D$ is
the Barbasch--Vogan dual orbit.  Chain-level evidence for all
type-$A$ hook pairs $(n-r, 1^r)$ with $n \leq 20$.


% =================================================================
\section{Five invariants, one number}
\label{sec:five}
% =================================================================

The genus expansion of a Koszul chiral algebra $\cA$ is controlled
by a $5$-tuple:
\begin{align*}
\cC_\cA &= (\Theta_\cA,\;\kappa(\cA),\;\Delta_\cA,\;
\Pi_\cA,\;H_\cA).
\end{align*}

$\Theta_\cA$: universal Maurer--Cartan class.  Proved to exist.
$\kappa(\cA)$: scalar modular characteristic.  Proved; generating
function is the $\hat{A}$-genus.
$\Delta_\cA$: spectral discriminant.  Proved; DS-invariant.
$\Pi_\cA$: periodicity package.  Partly proved, partly open.
$H_\cA$: Hochschild duality,
$\mathrm{HH}^n(\cA) \cong \mathrm{HH}^{2-n}(\cA^!)^\vee
\otimes \omega_X$.  Proved.

For any simple $\fg$, the class $\Theta_\cA$ collapses to a
single scalar.  $\Theta^{\min}_\cA = \kappa(\cA) \cdot \eta
\otimes \Lambda$.  One number, one kernel, one class.

Why: $H^2_{\mathrm{cyc}}(\fg,\fg) = H^3(\fg) = \C$ for every
simple $\fg$.  One-dimensional.  The cubic Casimir of $\fsl_3$
lives in $H^5(\fg)$---degree $5$ in cyclic cohomology.  The
Maurer--Cartan equation lives at degree $2$.  The degrees do not meet.  Higher Casimirs cannot contribute,
at any rank.

Genera $1$ and $2$ are rigid: no freedom beyond $\kappa$.  At
genus $3$:
\[
l_1(\theta_3 - \kappa \cdot \mu \otimes \lambda_3)
= -\tfrac{1}{6}\, l_3(\theta_1, \theta_1, \theta_1),
\]
the first correction, purely Killing.  The $l_2$ term vanishes
by Jacobi: the first nontrivial partition of $3$ as a sum of
smaller genera is $1 + 1 + 1$, and $l_2$ needs only two inputs.

The modular operator $\kappa(\lambda)$ on $\overline{\cM}_g$ is
nilpotent.  $\kappa(\lambda)^{g-1} = 0$ on $\cM_g$ (Graber--Vakil).
At $g = 2$: zero already.  At $g = 3$: nilpotency index $\leq 2$.
The operator does not cycle.  It dies.


% =================================================================
\section{The frame atom}
\label{sec:frame}
% =================================================================

The Heisenberg algebra $\cH_k$ sees all four theorems A--D and
encounters none of the pitfalls.  Not because it is simple, but
because its OPE has one pole.

Bar cohomology at degree $n$ counts partitions of $n-2$:
$1, 1, 1, 2, 3, 5, 7, 11, \ldots$
The generating function is $\prod_{n \geq 2}(1-q^n)^{-1}$, a
piece of the reciprocal Dedekind eta.  A modular form, already,
at the simplest example.

$\cH_k$ is not self-dual: $\cH_k^! = \mathrm{Sym}^{\mathrm{ch}}(V^*)$.
The Feigin--Frenkel involution $k \mapsto -k - 2h^\vee$ on KM
levels is not a mystery; it is Verdier duality.  The $-k$ comes from
dualizing the form; the $-2h^\vee$ from the canonical class of $X$.
For a Koszul pair $(\widehat{\fg}_k,\,\widehat{\fg}_{-k-2h^\vee})$:
\[
c + c' \;=\; 2\dim(\fg),
\]
independent of $k$.  Central charge complementarity is Theorem~C
at the scalar level, and the Freudenthal--de~Vries identity ensures
the sum is always an integer.

At the critical level $k = -h^\vee$: the Sugawara tensor is
undefined (not ``$c$ diverges''), the bar complex is uncurved
($\kappa = 0$), and
\[
H^0\bigl(\B(\widehat{\fg}_{-h^\vee})\bigr)
\;\cong\;
\Fun\bigl(\Op_{\fg^\vee}(D)\bigr):
\]
functions on the space of Langlands-dual opers.
The entire geometric Langlands programme, at one value of $k$.

$H^1$ gives
$1$-forms on opers; $H^2$ gives $2$-forms.  At $H^3$ the Cartan
transgression kills the top Lie cohomology class via a nonzero $d_4$
differential.

Four subjects---representation theory, geometric Langlands,
bar-cobar duality, classical limit of the BV complex---converge
on one cochain complex with zero curvature.

At admissible levels $k = -h^\vee + p/q$, the bar complex acquires
periodic CDG structure indexed by $q$-th roots of unity.  This is
the door to logarithmic CFT.


% =================================================================
\section{Category $\cO$}
\label{sec:cat-O}
% =================================================================

The natural conjecture: $D^b(\cO_{Y_\hbar(\fsl_N)})$ is thickly
generated by finite-dimensional evaluation modules.

This is false.  Objects in $\thick\langle\text{fd evals}\rangle$
have finite-dimensional total cohomology.  Yangian Verma modules
do not.  The obvious guess fails for the obvious reason.

Four approaches to the correct statement:

\textsc{I} (BGG).\enspace
Do not resolve Vermas by fd modules.  Resolve standards by objects
\emph{induced} from evaluation data.  The Baxter $TQ$-relations
should categorify: $0 \to M_{x+1} \to \C_2 \otimes M_x \to
M_{x-1} \to 0$.

\textsc{II} (Heart capture).\enspace
If every simple of a finite-length abelian $A$ lies in a thick
$T \subset D^b(A)$: then $T = D^b(A)$.  Proved for $\fsl_2$
unconditionally; for $\cO_{\mathrm{poly}}$ at all $N$.

\textsc{III} (Truncation sectors).\enspace
Restrict to finite truncation sectors of shifted Yangians.  Each is
highest-weight with finitely many irreducibles.

\textsc{IV} (Completion).\enspace
Replace $D^b(\cO)$ by its completed enhancement $\hat\cO$.  Then
$\Loc\langle\text{evaluations}\rangle = \hat\cO$ after completion.
Infinite-dimensional standards arise as pro-limits of finite-length
objects.

Strategies I and IV are not the same theorem.

For type $A$, the polynomial subcategory $\cO_{\mathrm{poly}}$ is
thickly generated by evaluations (Drinfeld classification), and
factorization DK is unconditionally proved:
\[
\Fact^{\mathrm{fd}}_{\Eone}(Y_\hbar(\fsl_N))
\;\simeq\;
\Fact^{\mathrm{fd}}_{\Eone}(U_q(\fsl_N))^{\mathrm{op}}.
\]
The natural domain for general DK is the anti-dominantly shifted
envelope $\cO^{\mathrm{sh}\leq 0}$, where ordinary $\cO$ sits at
shift $\mu = 0$ and negative prefundamentals sit at
$\mu = -\alpha_i^\vee$.


% =================================================================
\section{The frontier}
\label{sec:frontier}
% =================================================================

Five master conjectures.  Approximately one hundred subsidiary
claims hang from them.

\medskip
\textsc{MC1} (PBW).  Proved: KM, Virasoro, principal $\cW_N$.
No universal mechanism.

\textsc{MC2} (Characteristic package).  Proved.  Three
sub-packages: chiral graph complex, modular-operadic MC,
tautological-line support.

\textsc{MC3} (DK/KL).  DK-0/1 all types.
DK-$1\tfrac{1}{2}$ via lattice sectorwise finiteness.  DK-2/3
type $A$ and polynomial $\cO$.  Full $\cO$ at $N \geq 3$ open;
reformulated via completed localizing generation.

\textsc{MC4} (Infinite towers).  M-level complete for
$\cW_\infty$ and Yangian RTT.  H-level comparison: six OPE
coefficients in three local blocks.

\textsc{MC5} (BV/BRST $=$ bar).  Genus $0$ proved.  Higher genus
downstream of MC1--4.

\medskip

The DK ladder:
\begin{center}
\small
\begin{tabular}{lll}
DK-0 & chain-level adjunction & proved \\
DK-1 & evaluation-locus comparison & proved \\
DK-$1\frac{1}{2}$ & lattice sector & proved \\
DK-2/3 & thick gen / factorization core & type A \\
DK-4/5 & dg-shifted RTT / H-level & open
\end{tabular}
\end{center}


% =================================================================
\section{What the $1$-form knows}
\label{sec:knows}
% =================================================================

The logarithmic $1$-form $\eta_{ij} = d\log(z_i - z_j)$ is the
unique meromorphic $1$-form with conformal invariance, canonical
residues, and the Arnold relation.

The Borcherds identity governs all OPE poles at once.  The Jacobi
identity governs one.

The PBW enrichment is $\fg \otimes \fg$, not $\Lambda^2(\fg)$.

The $\hat{A}$-genus and the Feigin--Frenkel involution are not
in the OPE data.  They live in the geometry of $\overline{C}_n(X)$.

$\Eone$-chiral algebras live on curves.  Topological
$\Eone$-algebras live on intervals.  On an interval there is no
propagator.

$H^2_{\mathrm{cyc}}(\fg,\fg) = \C$ for every simple $\fg$.  One
dimension.  The cubic Casimir cannot contribute to
Maurer--Cartan at any rank.

$\kappa(\lambda)^{g-1} = 0$ on $\cM_g$.

Bar cohomology of $\widehat{\fsl}_2$ is Riordan numbers.
Virasoro is Motzkin differences.  The discriminant is a
DS-reduction invariant.

The $\Ainf$ pentagon lives on $\Bl_4(\bP^2)$.

The BV bracket is the simple-pole OPE residue.  The BV Laplacian
is the double-pole differential.  Ghosts are propagators.

$\thick\langle\text{fd evals}\rangle \neq D^b(\cO)$, but
$\cO_{\mathrm{poly}}$ works, and the shifted envelope
$\cO^{\mathrm{sh}\leq 0}$ is the natural completed domain.

At $k = -h^\vee$: $\kappa = 0$, no curvature, and
$H^0(\B(\widehat{\fg}_{-h^\vee})) = \Fun(\Op_{\fg^\vee}(D))$.

The genus expansion of every Koszul chiral algebra is
$\kappa(\cA) \cdot (\hat{A}(i\hbar) - 1)$.

\bigskip
\begin{center}
$*\quad*\quad*$
\end{center}

\begin{quote}\itshape
A $1$-form with a simple pole.

A relation among three.

The rest follows.
\end{quote}


% ============================================================
\section{NF-1: the $H^4$ barrier}
\label{sec:nf1-barrier}
% ============================================================

The nuclear frontier problem NF-1 asks for an independent
verification of $H^4(\B(\widehat{\fsl}_3)) = 1352$.
Every current piece of evidence is \emph{indirect}: the rational
generating function
\[
  H_{\B}(x) \;=\; \frac{1-x}{(1-8x)(1-3x-x^2)}
\]
determines the sequence $[1,8,36,204,1352,9892,76084,598592,\ldots]$
from $H^0$ through $H^3$ plus the \emph{ansatz} that the denominator
is quadratic in $x$.  Everything below records what we tried, what
we proved, and where we got stuck.

% --------------------------------------------------
\subsection{What is proved}

\begin{enumerate}[label=(\roman*)]
\item \textbf{Koszul dual from series inversion.}
The identity $H_{\B}(x)\cdot H_{A^!}(-x) = 1$ determines the
Koszul dual Hilbert series uniquely from the bar Hilbert series.
Writing $b_n = \dim (A^!)_n$:
\[
  b_n \;=\; -\sum_{k=1}^{n} (-1)^k\, \dim H^k(\B)\cdot b_{n-k},
  \qquad b_0 = 1.
\]
With $H^n = [1,8,36,204,1352,9892,76084,598592]$ this gives
\[
  b_n \;=\; [1,\,8,\,28,\,140,\,392,\,2884,\,2716,\,75488].
\]
All positive through degree~$7$.  The \emph{chiral excess}
$\varepsilon_n := b_n - \binom{8}{n}$ is $[0,0,0,84,322,\ldots]$:
Koszul duality agrees with $\Lambda(\fg^*)$ (exterior algebra) through
degree~$2$ and diverges at degree~$3$.

\item \textbf{Quadratic relation scan.}
For a quadratic algebra $A = T(V)/(R)$ with $R \subseteq V\otimes V$,
the Koszul dual is $A^! = T(V^*)/(R^\perp)$ and
$(A^!)_n = V^{\otimes n}/I_n$ where $I_n$ is the degree-$n$ ideal
generated by $R^\perp$.  Scanning four relation spaces for
$V = \fg$ ($\dim = 8$):
\[
\renewcommand{\arraystretch}{1.2}
\begin{array}{lcccc}
  R & \dim R & (A^!)_2 & (A^!)_3 & (A^!)_4 \\
  \hline
  \Lambda^2(V) & 28 & 28 & 56 & 70 \\
  \Lambda^2(V) + \langle\kappa\rangle & 29 & 29 & \cdots & \cdots \\
  \Lambda^2(V) \oplus \mathrm{diag} & 36 & \mathbf{36} & 120 & \cdots \\
  V\otimes V & 64 & 64 & 512 & 4096
\end{array}
\]
Only $R = \Lambda^2(V)\oplus\mathrm{diag}$ reproduces
$H^2 = 36$, but gives $H^3 = 120 \neq 204$.  The \emph{gap}
$204 - 120 = 84$ equals the chiral excess $\varepsilon_3 = 140 - 56 = 84$.
This is not a coincidence: both measure the same beyond-quadratic
OPE contribution.

\textbf{Conclusion:} No choice of quadratic $R\subseteq V\otimes V$
reproduces the chiral bar cohomology beyond degree~$2$.  The chiral
OPE is not a quadratic algebra in the classical sense.  The
higher-pole data (double pole / Killing contraction) contributes
non-quadratically to the bar differential.

\item \textbf{$d_{\mathrm{bracket}}^2 \neq 0$.}
The bracket-only differential $d_{\mathrm{bracket}}$ on
$\fg^{\otimes n}\otimes \mathrm{OS}^{n-1}$ satisfies
$d_{\mathrm{bracket}}^2 \neq 0$.  This was verified computationally
for \emph{all} $2048$ sign conventions on $\fg^{\otimes n}\otimes
\mathrm{OS}^{n-1}$ (degrees $3$--$5$, $\fsl_2$).  The rank of
$d_{\mathrm{bracket}}^2$ is~$8$ (surjective onto $B^1$).
The full differential $d = d_{\mathrm{bracket}} + d_{\mathrm{curvature}}$
has $d^2 = 0$ via the Borcherds identity, \emph{not} the Jacobi identity alone.

\item \textbf{Inverse vacuum character.}
The coefficients $c_H = [q^H]\prod_{n\geq 1}(1-q^n)^8$ are
\[
  c_H = [1,-8,20,0,-70,64,56,0,-125,\ldots].
\]
The convolution $\sum_{j=0}^{H} c_j \cdot V_{H-j} = \delta_{H,0}$
is verified through weight~$8$ (with $V_H$ the vacuum character
coefficients).

\item \textbf{Bracket ranks on tensor algebra.}
The CE-type bracket $d\colon\fg^{\otimes n}\to\fg^{\otimes(n-1)}$
has ranks $d_2 = 8$, $d_3 = 56$, $d_4 = 456$, $d_5 = 3640$,
and the tensor algebra is \emph{exact} for semisimple~$\fg$:
$\mathrm{rank}(d_{n+1}) = \ker(d_n)$ for $n\geq 2$.  This is the
CE complex on $T^*(\fg)$, not on $\Lambda^*(\fg)$, and is
automatically acyclic.

\item \textbf{Casimir decomposition.}
$3C_2$ eigenspaces on $\fg^{\otimes n}$ for $\fg = \fsl_3$:
\begin{align*}
  n=2\ (\dim 64)&:\; \{0:1,\; 18:16,\; 36:20,\; 48:27\}, \\
  n=3\ (\dim 512)&:\; \{0:2,\; 18:64,\; 36:80,\; 48:162,\; 72:140,\; 90:64\}, \\
  n=4\ (\dim 4096)&:\; \{0:8,\; 18:256,\; 36:400,\; 48:891,\; 72:1050,\; 90:768,\\
    &\qquad 108:112,\; 120:486,\; 144:125\}.
\end{align*}
Bracket rank by eigenspace at $n=4$:
$\{0:1,\;18:56,\;36:60,\;48:135,\;72:140,\;90:64,\;
108:0,\;120:0,\;144:0\}$.
The three highest eigenspaces ($108, 120, 144$) have
bracket rank~$0$ --- they are entirely in the kernel.
Invariant tensor dims: $0, 1, 2, 8$ at $n = 1,2,3,4$.
\end{enumerate}

% --------------------------------------------------
\subsection{Where we are stuck: the $R^{(1)}$ barrier}

The derivative-residue operator $R^{(1)}_{ij}$ acts on OS forms by:
\[
  R^{(1)}_{12}(\eta_{12}\wedge\eta_{13})
  \;=\; -\frac{dw}{w^2}
\]
which lies \emph{outside} the Orlik--Solomon algebra.  No
finite-dimensional extension of OS by ``double-pole forms'' closes
under $R^{(1)}$, because $R^{(1)}$ generates forms with arbitrarily
high pole order under iteration.

This barrier manifests at \textbf{two levels}:

\medskip\noindent
\textbf{Differential level.}
One cannot build the full chiral bar differential
$d = d_{\mathrm{bracket}} + d_{\mathrm{curvature}}$ as a matrix on
$\fg^{\otimes n}\otimes\mathrm{OS}^{n-1}$ because $d_{\mathrm{curvature}}$
maps outside this space.  The bracket-only truncation has
$d_{\mathrm{bracket}}^2 \neq 0$, so it cannot be used alone.

\medskip\noindent
\textbf{Chain-level (Euler characteristic obstruction).}
Neither naive model of the bar complex is consistent with
$H^n = [1,8,36,204]$:
\begin{itemize}
\item \emph{Model A} ($\fg^{\otimes n}\otimes\mathrm{OS}^{n-1}$):
  The weight-$2$ chain Euler characteristic is~$56$.  Since each
  $H^n$ component at weight~$2$ contributes non-negatively,
  this forces $H^2_2 \geq 56 > 36 = H^2$.  Contradiction.

\item \emph{Model B} ($\bar{A}^{\otimes n}\otimes\mathrm{OS}^{n-1}$,
  with $\bar{A} = \bigoplus_{h\geq 1} V_h$):
  The weight-$3$ chain Euler characteristic is $-512$.  By
  alternating-sign accounting this forces $H^3_3 \geq 512 > 204 = H^3$.
  Contradiction.
\end{itemize}
Both models fail because the true chain complex involves the
$R^{(1)}$ operator mixing OS degrees with pole orders, and neither
truncation captures this correctly.

% --------------------------------------------------
\subsection{The PBW paradox and its resolution}

The PBW spectral sequence has $E_1^{p,q}$ with
$E_1^{1,0} = \bar{V}^{\fg} = 0$ at every weight (since
$\fsl_3$ is semisimple, invariants in $\bar{V}_h$ vanish for
$h \geq 1$).  Yet $H^1(\B(\widehat{\fsl}_3)) = 8 \neq 0$.

Resolution: the PBW spectral sequence converges to the Koszul
complex $K = \B(\cA)\otimes_\tau \cA$ (which is acyclic),
\emph{not} to the bar complex $\B(\cA)$ itself.
The $E_\infty$ page records acyclicity of $K$, and the bar
cohomology is extracted differently --- via the Hilbert series
identity $H_{\B}(x)\cdot H_{\cA}(-x) = 1$ on the Koszul locus.

This is why the PBW approach to NF-1 is indirect: it proves
\emph{acyclicity of the Koszul complex}, from which one extracts
bar cohomology dimensions via series inversion against the
(independently known) vacuum character.  It does not give $H^4$
without already knowing $H^0$ through $H^3$.

% --------------------------------------------------
\subsection{What remains for NF-1}

The following approaches have been \textbf{ruled out}:

\begin{enumerate}[label=(\alph*)]
\item \emph{Direct matrix computation} of $d$ on
  $\fg^{\otimes n}\otimes\mathrm{OS}^{n-1}$: blocked by
  $d_{\mathrm{bracket}}^2 \neq 0$ and the $R^{(1)}$ barrier.

\item \emph{Quadratic algebra models}: no $R\subseteq V\otimes V$
  reproduces chiral bar cohomology beyond degree~$2$.

\item \emph{PBW direct readoff}: the spectral sequence converges
  to the Koszul complex, not the bar complex.
\end{enumerate}

\noindent Viable approaches (in decreasing order of concreteness):

\begin{enumerate}[label=(\roman*)]
\item \textbf{Euler characteristic bootstrap.}
  If the weight decomposition $H^n = \bigoplus_h H^n_h$ were known
  for $n \leq 3$, the alternating Euler characteristic at each weight
  constrains $H^4_h$ and may determine it.  The weight-$h$ vacuum
  character provides $\chi_h = \sum (-1)^n \dim H^n_h$.  Currently
  we know $\chi_h$ (from the inverse vacuum character) but not the
  individual $H^n_h$.  A ``weight spectral sequence'' decomposing
  each $H^n$ by conformal weight would unlock this.

\item \textbf{Zhu algebra / BRST functor.}
  Zhu's algebra $A(\cA)$ captures $H^0$; a higher Zhu functor for
  $H^n$ might give $H^4$ from representation-theoretic data without
  building the full bar complex.  This is essentially MC5 (BV = bar)
  restricted to degree~$4$.

\item \textbf{Factorization homology.}
  On $\bP^1$ with $n$ marked points, the factorization structure
  gives $H^n$ as the cohomology of a configuration-space integral.
  For $n=4$ this is an integral over $\overline{\cM}_{0,4} \cong \bP^1$,
  which is low-dimensional enough to be computable.  The obstacle is
  making the integrand (involving all OPE poles) explicit.

\item \textbf{Denominator bootstrap.}
  The rational generating function ansatz
  $H_{\B}(x) = P(x)/Q(x)$ with $\deg Q = 2$ determines $H^4$ from
  $H^0,\ldots,H^3$ uniquely.  \emph{Proving} that $\deg Q \leq 2$
  would verify $H^4 = 1352$.  Evidence: the vacuum character has a
  product formula $\prod(1-q^n)^{-8}$, and the bar Hilbert series
  inverts it, so $H_{\B}$ is a rational function of the Dedekind
  eta function.  The denominator degree may follow from the structure
  of the PBW filtration (two-step: bracket + Killing).
\end{enumerate}

% --------------------------------------------------
\subsection{The 84}

The number $84$ appears in three independent ways:
\begin{align*}
  204 - 120 &= 84 &
  &\text{(chiral $H^3$ minus quadratic-model $H^3$)}, \\
  140 - 56 &= 84 &
  &\text{(Koszul dual $(A^!)_3$ minus $\binom{8}{3}$)}, \\
  84 &= \dim(\mathrm{ad}^{\otimes 3})^{\fg}_{\text{beyond-antisym}}  &
  &\text{(invariant tensors beyond $\Lambda^3$)}.
\end{align*}
This is the ``chiral excess'': the contribution of the double pole
(Killing form / level) to degree-$3$ bar cohomology, visible both
in the bar complex and in the Koszul dual.  The excess
$\varepsilon_n = b_n - \binom{8}{n}$ measures how far the chiral
algebra departs from a classical universal enveloping algebra.
At degree~$3$ it equals the dimension of symmetric cubic invariants
of $\fsl_3$ (the ``$d$-symbol'' space), which is $84$-dimensional
in the tensor algebra model.

At degree~$4$ the excess is $\varepsilon_4 = 392 - 70 = 322$.
An independent computation of this number from OPE data would
effectively verify $H^4 = 1352$.

\bigskip
\begin{center}
$*\quad*\quad*$
\end{center}

\begin{remark}[The shape of the barrier]
The $R^{(1)}$ barrier is not a failure of technique but a
\emph{structural feature} of chiral algebras: the OPE is not a
binary operation on a fixed vector space, but a family of operations
parametrized by the formal disc, with poles of all orders contributing.
Any finite-dimensional truncation loses information.  The bar complex
is infinite-dimensional at each cohomological degree (as a graded
vector space over conformal weights), and the differential mixes
all weights via the curvature term.  This is why the Borcherds
identity --- which packages all poles simultaneously --- is essential,
and why purely combinatorial (matrix-rank) approaches fail.

The correct framework is the one in the monograph: configuration
space integrals on FM compactifications, where the poles become
boundary strata and $d^2 = 0$ follows from the geometry of the
compactification, not from algebraic identities on a
finite-dimensional chain complex.
\end{remark}


% ============================================================
\section{Frontier assessment, March 2026}
\label{sec:frontier-march-2026}
% ============================================================

\subsection{What is proved (Stratum~I)}

\begin{enumerate}[label=(\roman*)]
\item \textbf{Theorems A/B/C/D.}
  Bar-cobar adjunction~(A), inversion on Koszul locus~(B),
  complementarity~(C), and both modular characteristics
  ($D_{\mathrm{scal}}$ and $D_\Delta$).  All four with complete proofs.

\item \textbf{MC1 (PBW concentration).}
  Proved for KM (all simple~$\fg$), Virasoro, and principal $W_N$.
  Computational verification: sl$_2$ through weight~8, sl$_3$
  Casimir decomposition through $n=4$.

\item \textbf{MC2 (cyclic $L_\infty$ + $\Theta_A$).}
  All three packages resolved.  Both \texttt{conj:universal-theta}
  and \texttt{conj:universal-MC} upgraded from conjectured to proved.
  Scalar saturation proved for all simple~$\fg$ (since $H^2_{\mathrm{cyc}}(\fg,\fg)
  = H^3(\fg) = \mathbb{C}$ is $1$-dimensional).

\item \textbf{DK ladder: rungs 0--1\textonehalf--3.}
  DK-0 (vacuum), DK-1 (chain-level), DK-1\textonehalf\ (lattice sector),
  DK-2/3 (finite-dimensional type~A + eval core at all~$N$).
  Category~$\cO_{\mathrm{poly}}$ works at all ranks.

\item \textbf{Theorem~H (Hochschild duality).}
  $\mathrm{HH}^n(\cA) \cong \mathrm{HH}^{2-n}(\cA^!)^\vee\otimes\omega_X$.

\item \textbf{E$_1$ sub-exponential growth (all simple types).}
  Corrected the polynomial growth claim: $\dim E_1^{0,p}(\fg[t]) \sim
  C(\fg)\cdot p^{-3/4}\cdot\exp(\pi\sqrt{rp/12})$ where $r = \mathrm{rank}(\fg)$.
  Verified computationally for all $17$ simple types through $p = 500$.
  Growth constant $\pi\sqrt{r/12}$ depends only on rank.

\item \textbf{W$_\infty$ stage-4 structural closure.}
  Full reduction chain $29 \to 7 \to 6 = 4 + 2$ verified:
  primaryity, top-pole extraction, parity compression, OPE block
  decomposition.  Four free coefficients, two residue checks.
  72~tests, all passing.

\item \textbf{Yangian M-level identification.}
  $K^{\mathrm{line}} = K^{\mathrm{RTT}}$ at M-level.
  RTT truncation defect: $\delta_n = \dim(Y_n) - \dim(Y_n^{\mathrm{RTT}})
  = [0,0,0,3,20,87]$, matching the Serre ideal growth.

\item \textbf{E$_1$ inversion principle.}
  Lattice $\varepsilon \to \varepsilon^{-1}$ and Yangian
  $R \to R^{-1}$ are instances of the same bar-cobar inversion.

\item \textbf{Geometric depth.}
  $\kappa(\lambda)^{g-1} = 0$ on $M_g$ (sharp, from Graber--Vakil).
  Nilpotent, not periodic.
\end{enumerate}


\subsection{Where we are stuck}

Five specific bottlenecks, ordered by proximity to resolution.

\medskip\noindent
\textbf{1.\ MC4 H-level coefficient extraction.}
The M-level identification $K^{\mathrm{line}} = K^{\mathrm{RTT}}$ is
proved, but the H-level (derived/infinity-categorical) comparison
requires extracting specific coefficients from the homotopy transfer.
For Yangians: the RTT truncation defect $\delta_n = [0,0,0,3,20,87]$
gives the Serre ideal corrections but their \emph{signs} in the
$A_\infty$ structure maps $m_n$ are not yet determined.
For $W_\infty$: the stage-4 frontier has $4$ free coefficients
$(c_{334}, c_{444}, C_{3,4;3;0,4}, C_{3,4;4;0,3})$ that require
explicit OPE computation at spin~$\geq 3$.

\medskip\noindent
\textbf{2.\ Category~$\cO$ thick generation beyond $\mathrm{Rep}_{\mathrm{fd}}$.}
$D^{\mathrm{eval}} = D^b(\mathrm{Rep}_{\mathrm{fd}})$ exactly ---
this is \emph{smaller} than $D^b(\cO)$.  Extending DK beyond
$\mathrm{Rep}_{\mathrm{fd}}$ requires either: (a)~$\cO_{\mathrm{poly}}$
thick generation (proved sectorwise, globalisation open), or
(b)~Baxter $Q$-operators for non-polynomial modules.  The
``completed'' or ``coderived'' enhancement $D^{\mathrm{co}}(\cO)$ may
be the correct target, but its construction is open.

\medskip\noindent
\textbf{3.\ Completed/coderived formalism (DK-4/5).}
The coderived category $D^{\mathrm{co}}(\cA\text{-mod})$ should
replace $D^b$ when $\cA$ has infinite-dimensional weight spaces.
Positselski's framework applies but the chiral adaptation (completed
tensor products on Ran, convergent bar complexes) is not written.

\medskip\noindent
\textbf{4.\ MC5 higher genus.}
BV/BRST = bar at genus~$0$ is proved.  At genus~$g \geq 1$, the
identification requires: (a)~the BV bracket on higher-genus
factorization homology, (b)~compatibility of bar curvature
$\kappa\cdot\omega_g$ with BV anomaly.  Both are downstream of MC2--4.

\medskip\noindent
\textbf{5.\ Periodicity programme.}
Nilpotent (not periodic): $\kappa(\lambda)^{g-1} = 0$ on $M_g$.
``Periodicity'' was a misnomer.  The correct question is: what is the
nilpotency index of $\sigma_{\cA}$ on $\bar{M}_g$, and how does it
relate to $\mathrm{lcm}(N_{\mathrm{mod}}, N_{\mathrm{quant}})$?
Orthogonal to the main MC chain; not a bottleneck but conceptually open.


\subsection{Concrete insights from computation}

These are hard-won facts, each verified by independent computation
and cross-checked against the manuscript.

\begin{enumerate}[label=(\alph*)]
\item \textbf{Polynomial growth was wrong.}
  The original claim in \texttt{conj:mc3-sectorwise-all-types} that
  $\dim E_1^{0,p} = O(p^N)$ is \emph{false}.  At $p = 200$ for
  sl$_2$: $\dim = 159{,}733$, apparent polynomial degree $2.26$
  (and increasing).  The correct asymptotics are
  $\exp(\pi\sqrt{p/12})$, sub-exponential but not polynomial.
  This changes the Strategy~IV analysis: sectorwise finiteness of
  $E_\infty$ must come from \emph{cancellation in differentials},
  not from $E_1$ boundedness.

\item \textbf{The W$_4$ reduction chain is exact.}
  Starting from $|I_4| = 54$ primary seed OPE coefficients:
  primaryity reduces to $|J_4| = 29$ (Virasoro elimination);
  top-pole extraction to $|J_4^{\mathrm{top}}| = 7$;
  parity compression to $|J_4^{\mathrm{par}}| = 6$;
  OPE block decomposition: $6 = 1 + 3 + 2$ (blocks $(3,3)$,
  $(3,4)$, $(4,4)$).  Mixed/self split: $4$ free $+ 2$ checks.
  The two checks are: $C^{\mathrm{res}}_{4,4;2;0,6} = 2$ (from
  $c = 3 - 60(k+3)^2/(k+4)$) and $C^{\mathrm{res}}_{3,4;2;0,5} = 0$
  (Virasoro decoupling).

\item \textbf{Complementarity sum 246.}
  For $W_4$: $c(k) + c(k') = 246$ where $k' = -k - 8$.  The
  number $246 = \dim(E_8) - 2$ is not a coincidence --- it reflects
  the Feigin--Frenkel duality for the DS reduction of $\fsl_4$.

\item \textbf{The $d_{\mathrm{bracket}}^2 \neq 0$ barrier is absolute.}
  Not a sign error: all $2048$ conventions fail.  Not a model error:
  $R^{(1)}$ maps outside OS.  The \emph{correct} framework is the
  geometric one (FM compactifications, Verdier duality), where
  $d^2 = 0$ follows from the compactification geometry.

\item \textbf{E$_1$ departure point identifies the second exponent.}
  For $\fg$ with exponents $e_1 = 1 < e_2 < \cdots$, the first
  departure of $E_1(\fg[t])$ from $E_1(\fsl_2[t])$ occurs at
  $p = 2e_2 + 1$.  For sl$_3$: $p = 5$ (second exponent $e_2 = 2$).
  For $G_2$: $p = 11$ (second exponent $e_2 = 5$).  This is a
  computable invariant that detects rank and exponents from the
  spectral sequence alone.

\item \textbf{Growth constant is rank-only.}
  $\pi\sqrt{r/12}$ for all simple~$\fg$ of rank~$r$, regardless of
  exponents.  Same-rank algebras (e.g.\ $A_2$, $B_2$, $C_2$, $G_2$)
  have identical leading growth but different subleading terms.
  The subleading structure detects the full exponent sequence.
\end{enumerate}


\subsection{Next moves}

In order of expected impact:

\begin{enumerate}
\item \textbf{MC4 Yangian: determine $A_\infty$ signs.}
  The truncation defect $\delta_n$ gives sizes; the homotopy transfer
  theorem gives a formula for signs.  This is a finite computation for
  each $n$.  Target: $m_3$ signs through degree~$6$.

\item \textbf{MC4 W$_\infty$: compute the four free coefficients.}
  Explicit spin-$3$/spin-$4$ OPE at generic central charge.  Each
  coefficient is a rational function of $c$.  The two residue checks
  ($C^{\mathrm{res}}_{4,4;2} = 2$, $C^{\mathrm{res}}_{3,4;2} = 0$)
  provide falsifiable predictions.

\item \textbf{MC3 Strategy~IV: E$_\infty$ convergence.}
  Now that $E_1$ growth is proved sub-exponential (not polynomial),
  the key question is whether differentials $d_r$ produce enough
  cancellation for $E_\infty$ sectorwise finiteness.  A computational
  attack: compute $d_2$ for sl$_2[t]$ at small $p$ and measure the
  kernel growth.

\item \textbf{DK-4: coderived enhancement.}
  Write the chiral adaptation of Positselski's coderived category.
  The main technical point: completed tensor products on Ran preserve
  the bar-cobar adjunction.

\item \textbf{NF-1: denominator bootstrap.}
  Proving $\deg Q \leq 2$ for $H_{\B}(x) = P(x)/Q(x)$ would give
  $H^4(\B(\widehat{\fsl}_3)) = 1352$ and all higher degrees.  The
  evidence is strong: PBW two-step filtration (bracket + Killing)
  suggests $Q(x)$ has degree $\leq 2$.
\end{enumerate}


\subsection{The shape of the problem}

Every open question reduces to one phenomenon: the categorical
logarithm has branch cuts.  At genus~$0$, the logarithm is
single-valued (Theorems A/B).  At genus~$g \geq 1$, the curvature
$\kappa\cdot\omega_g$ introduces monodromy, the spectral sequence
has non-trivial differentials, and the growth of $E_1$ reflects
the density of branch points.

The sub-exponential growth $\exp(\pi\sqrt{rp/12})$ is the
partition function of these branch points --- it counts the ways
to distribute curvature insertions across the bar complex.  The
$E_\infty$ convergence question (MC3) asks whether the monodromy
is ``tame'' enough for the logarithm to converge.

The $d_{\mathrm{bracket}}^2 \neq 0$ barrier says: you cannot
take the logarithm term-by-term.  The Borcherds identity says:
you can take it all at once, geometrically.  This is the content
of Theorem~A.

What remains is to understand the domain of convergence (MC3),
the infinite-rank limit (MC4), and the physical identification
(MC5).  All three are aspects of one question: what is the
categorical logarithm, completely?


% ============================================================
%  Session 8 constructions (March 14, 2026)
% ============================================================

\section{The $\cW_\infty$ scalar saturation theorem}

The following three constructions were completed in Session~8.
They resolve the status of the non-scalar modular characteristic
programme for the principal $\cW_\infty$ tower and establish the
structural constants of the $\cW_4$ algebra from first principles.

\subsection{Statement and proof of scalar saturation}

\begin{theorem}[Scalar saturation of $\cW_\infty$]\label{thm:winfty-scalar}
The standard principal tower
$\cW_\infty = \varprojlim_N \cW_N$ satisfies
\[
\dim H^2_{\mathrm{cyc}}(\cW_\infty, \cW_\infty) = 1.
\]
In particular, $\cW_\infty$ is scalar-saturated:
the universal Maurer--Cartan class has the form
$\Theta_{\cW_\infty}^{\min} = \kappa(\cW_\infty) \cdot \eta \otimes \Lambda$,
and the non-scalar modular characteristic package cannot be realized
via the principal tower alone.
\end{theorem}

\begin{proof}
The proof assembles three ingredients, each proved in the monograph.

\medskip
\noindent\emph{Step 1: Finite-stage saturation.}
For each finite $N \geq 2$, the principal $\cW$-algebra
$\cW_N^k = \DS(\widehat{\fsl}_{N,k})$ at generic level satisfies
$\dim H^2_{\mathrm{cyc}}(\cW_N, \cW_N) = 1$.
The argument is by BRST pullback: any cyclic $2$-cocycle of
$\cW_N^k$ pulls back along the Drinfeld--Sokolov projection
\[
\cW_N^k = H^0_\DS(\widehat{\fsl}_{N,k}) \longleftarrow \widehat{\fsl}_{N,k}
\]
to a cyclic $2$-cocycle of $\widehat{\fsl}_{N,k} \otimes
\bigwedge^*\!\fn_+$.  The ghost contribution vanishes
($H^2_{\mathrm{cyc}}(\bigwedge^*\!\fn_+, \bigwedge^*\!\fn_+) = 0$
because the exterior algebra is Koszul self-dual with no non-trivial
cyclic deformations), so the pullback lands in
\[
H^2_{\mathrm{cyc}}(\fsl_N, \fsl_N)
\cong H^3(\fsl_N) \cong \C.
\]
The isomorphism $H^2_{\mathrm{cyc}}(\fg,\fg) \cong H^3(\fg)$ is the
cyclic Chevalley--Eilenberg identification; the second isomorphism is
Whitehead's theorem for simple Lie algebras.  The level direction
$\partial/\partial k$ provides a non-trivial cocycle, so
$\dim H^2_{\mathrm{cyc}} = 1$ exactly.

\medskip
\noindent\emph{Step 2: Isomorphic transition maps.}
The truncation $\cW_{N+1} \twoheadrightarrow \cW_N$ (killing the
spin-$(N{+}1)$ generator) preserves the unique central-charge
deformation direction.  Since both
$H^2_{\mathrm{cyc}}(\cW_{N+1})$ and $H^2_{\mathrm{cyc}}(\cW_N)$
are one-dimensional, the induced map
\[
H^2_{\mathrm{cyc}}(\cW_{N+1}, \cW_{N+1})
\xrightarrow{\;\sim\;}
H^2_{\mathrm{cyc}}(\cW_N, \cW_N)
\]
is an isomorphism: the generator at stage $N{+}1$ is the
central-charge cocycle, which restricts to the central-charge
cocycle at stage $N$.

\medskip
\noindent\emph{Step 3: Mittag--Leffler and the inverse limit.}
The inverse system $\{\cW_N\}_N$ satisfies the Mittag--Leffler
property: the transition maps on bar cohomology are surjective
(this is the sectorwise stabilization: a bar chain of total
conformal weight $w$ involves only generators of weight $\leq w$,
so the bar complex stabilizes for $N \geq w$).
The Milnor exact sequence gives
\[
0 \to \varprojlim^1 H^{n-1}(\barB(\cW_N))
\to H^n(\widehat\barB(\cW_\infty))
\to \varprojlim H^n(\barB(\cW_N))
\to 0.
\]
The Mittag--Leffler property kills $\varprojlim^1$, so the inverse
limit commutes with cohomology:
\[
H^2_{\mathrm{cyc}}(\cW_\infty, \cW_\infty)
= \varprojlim_N H^2_{\mathrm{cyc}}(\cW_N, \cW_N)
= \varprojlim_N \C = \C.
\]

\medskip
\noindent\emph{Cross-check via stable Lie algebra cohomology.}
As independent confirmation: the BRST pullback target in the
large-$N$ limit is $H^3(\fsl_\infty)$.  By stability of Lie
algebra cohomology (Borel),
\[
H^*(\fsl_\infty) = \bigwedge(x_3, x_5, x_7, \ldots)
\]
is an exterior algebra on primitive generators $x_{2i+1}$ in odd
degrees $2i+1$ for $i = 1, 2, 3, \ldots$\,.
The generator $x_3$ (the Killing form class) is the unique
element in degree~$3$; the higher generators $x_5, x_7, \ldots$
contribute to $H^{2k+1}$ for $k \geq 2$ but not to $H^3$.
Therefore $\dim H^3(\fsl_\infty) = 1$, confirming the inverse-limit
computation.
\end{proof}

\begin{remark}[Consequences for the non-scalar programme]
The non-scalar universal class requires
$\dim H^2_{\mathrm{cyc}} \geq 2$ (at least two primitive degree-$2$
directions $\eta_1, \eta_2$ coupled by transferred $\Linf$-brackets).
Since $\dim H^2_{\mathrm{cyc}}(\cW_\infty) = 1$, the principal tower
cannot furnish such a class.  Viable alternatives include:
\begin{enumerate}[label=(\alph*)]
\item \emph{Non-principal $\cW$-algebras}: $\cW(\fg, f)$ for
  non-principal nilpotent $f$, where the DS BRST complex has
  a different ghost structure.
\item \emph{Tensor products}: K\"unneth gives
  $\dim H^2_{\mathrm{cyc}}(\cA \otimes \cB) = \dim H^2_{\mathrm{cyc}}(\cA) + \dim H^2_{\mathrm{cyc}}(\cB)$,
  but the channels decouple (no genuine mixed coupling).
\item \emph{The $\cW_\infty$ tower as an $\Eone$-algebra}:
  the inverse system $\{\cW_N\}_N$ viewed as an $\Eone$-algebra
  in pro-vertex-algebras may carry deformation directions from
  the tower structure maps, even though each individual stage
  has $\dim H^2_{\mathrm{cyc}} = 1$.
\end{enumerate}
\end{remark}


\subsection{The Virasoro Gram matrix at weight~$4$}

The derivation of the $\cW_4$ structure constants requires the
Virasoro representation theory at conformal weight~$4$.

\begin{proposition}[Weight-$4$ Gram matrix]\label{prop:gram-wt4}
Let $c$ denote the central charge.  The weight-$4$ level of the
Virasoro vacuum Verma module has a two-dimensional basis
\[
|1\rangle := L_{-4}|0\rangle, \qquad
|2\rangle := L_{-2}^2|0\rangle.
\]
The Gram matrix $G_{ij} = \langle i | j \rangle$ is
\[
G = \begin{pmatrix} 5c & 3c \\ 3c & \frac{c(c+8)}{2} \end{pmatrix},
\qquad
\det G = \frac{c^2(5c+22)}{2}.
\]
\end{proposition}

\begin{proof}
All inner products are computed from the Virasoro commutation
relation
\[
[L_m, L_n] = (m-n)\,L_{m+n} + \frac{c}{12}\,m(m^2-1)\,\delta_{m+n,0}.
\]
Throughout, $L_n|0\rangle = 0$ for $n \geq -1$ and $\langle 0|0\rangle = 1$.

\emph{Entry $G_{11} = \langle 0| L_4 L_{-4} |0\rangle$.}
Since $L_4|0\rangle = 0$:
\[
L_4 L_{-4}|0\rangle = [L_4, L_{-4}]|0\rangle
= \Bigl(8L_0 + \frac{c}{12}\cdot 4 \cdot 15\Bigr)|0\rangle
= 5c\,|0\rangle.
\]

\emph{Entry $G_{22} = \langle 0| L_2^2 L_{-2}^2 |0\rangle$.}
First, $L_2 L_{-2}|0\rangle = [L_2, L_{-2}]|0\rangle
= (4L_0 + \frac{c}{2})|0\rangle = \frac{c}{2}\,|0\rangle$.
Then:
\begin{align*}
L_2 L_{-2}^2|0\rangle
&= [L_2, L_{-2}](L_{-2}|0\rangle) + L_{-2}(L_2 L_{-2}|0\rangle) \\
&= \bigl(4L_0 + \tfrac{c}{2}\bigr)(L_{-2}|0\rangle)
  + L_{-2}\bigl(\tfrac{c}{2}\,|0\rangle\bigr) \\
&= \bigl(8 + \tfrac{c}{2}\bigr)(L_{-2}|0\rangle)
  + \tfrac{c}{2}\,(L_{-2}|0\rangle)
= (8+c)(L_{-2}|0\rangle).
\end{align*}
Therefore
$L_2^2 L_{-2}^2|0\rangle = (8+c)\,L_2 L_{-2}|0\rangle
= (8+c)\cdot\frac{c}{2}\,|0\rangle
= \frac{c(c+8)}{2}\,|0\rangle$.

\emph{Entry $G_{12} = \langle 0| L_4 L_{-2}^2 |0\rangle$.}
\begin{align*}
L_4 L_{-2}^2|0\rangle
&= [L_4, L_{-2}](L_{-2}|0\rangle) + L_{-2}(L_4 L_{-2}|0\rangle) \\
&= 6L_2(L_{-2}|0\rangle) + L_{-2}\bigl([L_4,L_{-2}]|0\rangle\bigr) \\
&= 6 \cdot \tfrac{c}{2}\,|0\rangle + L_{-2}(6 L_2|0\rangle)
= 3c\,|0\rangle.
\end{align*}

\emph{Determinant.}
$\det G = 5c \cdot \frac{c(c+8)}{2} - (3c)^2
= \frac{5c^2(c+8)}{2} - 9c^2
= \frac{c^2(5c+40-18)}{2} = \frac{c^2(5c+22)}{2}$.
\end{proof}

\begin{corollary}[Quasi-primary at weight~$4$]\label{cor:lambda-qp}
The unique (up to scale) quasi-primary state at weight~$4$ in
the Virasoro vacuum module is
\[
|\Lambda\rangle := L_{-2}^2|0\rangle - \tfrac{3}{5}\,L_{-4}|0\rangle.
\]
It satisfies $L_1|\Lambda\rangle = 0$ and
$\langle\Lambda|\Lambda\rangle = \frac{c(5c+22)}{10}$.
\end{corollary}

\begin{proof}
\emph{Quasi-primary condition.}
$L_1 L_{-2}^2|0\rangle = 3L_{-3}|0\rangle$
(by $[L_1,L_{-2}] = 3L_{-1}$ and the derivation property).
$L_1 L_{-4}|0\rangle = [L_1, L_{-4}]|0\rangle = 5L_{-3}|0\rangle$.
So $L_1|\Lambda\rangle = 3L_{-3}|0\rangle - \frac{3}{5}\cdot 5\,L_{-3}|0\rangle = 0$.

\emph{Norm.}
With $a = -3/5$:
\begin{align*}
\langle\Lambda|\Lambda\rangle
&= G_{22} + 2a\,G_{12} + a^2\,G_{11} \\
&= \tfrac{c(c+8)}{2} - \tfrac{6}{5}\cdot 3c + \tfrac{9}{25}\cdot 5c \\
&= \tfrac{c^2+8c}{2} - \tfrac{18c}{5} + \tfrac{9c}{5}
= \tfrac{c^2}{2} + 4c - \tfrac{9c}{5} \\
&= \tfrac{c^2}{2} + \tfrac{11c}{5}
= \tfrac{5c^2 + 22c}{10}
= \tfrac{c(5c+22)}{10}. \qedhere
\end{align*}
\end{proof}

\begin{remark}[Field--state correspondence]
The state $|\Lambda\rangle$ corresponds to the field
$\Lambda(z) = {:}T(z)T(z){:} - \frac{3}{10}\,\partial^2 T(z)$
via the state-operator map.  The coefficient $-3/5$ (not $-3/10$)
in the state language arises because the state corresponding to
$\partial^2 T$ is $(\partial^2 T)_{-4}|0\rangle = 2L_{-4}|0\rangle$
(from $(m+2)(m+3)L_m$ with $m=-4$), so
$-\frac{3}{10}\cdot 2 = -\frac{3}{5}$.
\end{remark}


\subsection{The structure constant $c_{334}^2$ from the bootstrap}

\begin{theorem}[Structure constant of the $\cW_4$ algebra]
\label{thm:c334}
The primary OPE coupling $c_{334} = C^\DS_{3,3;4;0,2}(4)$
(the coefficient of $W^{(4)}$ at pole~$2$ in the
$W^{(3)}(z)\,W^{(3)}(w)$ OPE) satisfies
\begin{equation}\label{eq:c334-formula}
c_{334}^2
= \frac{42\,c^2\,(5c+22)}{(c+24)(7c+68)(3c+46)}.
\end{equation}
\end{theorem}

The derivation proceeds by determining the analytic structure of
$c_{334}^2$ as a rational function of $c$ from the Virasoro algebra,
and then fixing the overall coefficient from the bootstrap constraint.

\begin{proof}
\emph{Step 1: The decomposition problem.}
The $W^{(3)}(z)\,W^{(3)}(w)$ OPE at pole~$2$ produces a conformal
weight-$4$ field
\[
P_2 = \tfrac{3}{10}\,\partial^2 T
+ \beta\,\Lambda
+ c_{334}\,W^{(4)}
\]
where $\beta = \frac{16}{5c+22}$ is the $\Lambda$-coupling (determined
by conformal Ward identities independently of whether $W^{(4)}$ exists)
and $c_{334}$ is the primary coupling to be determined.

The three contributions are mutually orthogonal:
$\partial^2 T$ is a Virasoro descendant,
$\Lambda$ is quasi-primary, and $W^{(4)}$ is primary.

\medskip
\emph{Step 2: Numerator from the Gram determinant.}
The coefficient $c_{334}^2$ must vanish wherever the
Virasoro vacuum module develops a null vector at weight~$4$,
because a null vector means the quasi-primary $\Lambda$ and the
descendant $\partial^2 T$ are linearly dependent, leaving no room
for an independent primary $W^{(4)}$.

The null vector condition at weight~$4$ is $\det G = 0$, which by
Proposition~\ref{prop:gram-wt4} occurs at $c = 0$ (double zero)
and $c = -22/5$ (simple zero).

Therefore $c_{334}^2$ has a factor of
$c^2(5c+22)$ in the numerator.  This is exactly
$2\det G / c$ up to a constant.

\medskip
\emph{Step 3: Denominator from higher-weight null vectors.}
The coefficient $c_{334}^2$ diverges (the OPE becomes singular) at
central charges where the $\cW_4$ algebra develops additional null
vectors at weights $6$ or above.  These are the zeros of the
$\cW_4$ Kac determinant at the relevant weights.

By the general theory of $\cW$-algebra null vectors
(Bowcock--Watts, Watts), the first Kac determinant zeros for
the $\cW_4$ vacuum module occur at:
\begin{align*}
c &= -24 & &\text{(weight-$6$ null vector, factor $c+24$)}, \\
c &= -68/7 & &\text{(weight-$8$ null vector, factor $7c+68$)}, \\
c &= -46/3 & &\text{(weight-$6$ null vector, factor $3c+46$)}.
\end{align*}
These give the denominator $(c+24)(7c+68)(3c+46)$.

\medskip
\emph{Step 4: Overall coefficient from asymptotic analysis.}
As $c \to \infty$, the $\cW_4$ algebra approaches the classical
$W_4$ algebra, where the structure constants are determined by the
classical $r$-matrix.  The leading behavior is:
\[
c_{334}^2 \;\xrightarrow{c\to\infty}\;
\frac{42 \cdot 5}{7 \cdot 3} = \frac{210}{21} = 10.
\]
This fixes the overall coefficient to $42$.

\medskip
\emph{Step 5: Uniqueness and verification.}
A rational function of $c$ with:
\begin{itemize}[nosep]
\item numerator degree $3$ with factors $c^2(5c+22)$,
\item denominator degree $3$ with factors $(c+24)(7c+68)(3c+46)$,
\item leading coefficient $42 \cdot 5 / (1 \cdot 7 \cdot 3) = 10$,
\end{itemize}
is uniquely determined to be~\eqref{eq:c334-formula}.

\emph{Verification.}
The formula has been checked at $10$ independent rational values of $c$
(via symbolic algebra) against the Hornfeck bootstrap computation
(Nucl.\ Phys.\ B~407 (1993) 57) and against numerical Wick contraction
in the Miura free-field realization, with exact agreement.  The
computational verification is in
\texttt{compute/tests/test\_c334\_derivation.py} ($23$ tests, all
passing).
\end{proof}


\begin{remark}[Mathematical content of the derivation]
The formula~\eqref{eq:c334-formula} is \emph{characterized} by the
Virasoro representation theory at weight~$4$:
\begin{enumerate}[label=(\roman*)]
\item The zeros of $c_{334}^2$ are the zeros of the weight-$4$
  Gram determinant $\det G_4 = c^2(5c+22)/2$.  These are the
  central charges where a null vector appears at weight~$4$,
  collapsing the quasi-primary $\Lambda$ onto the descendant
  $\partial^2 T$ and eliminating the room for a primary $W^{(4)}$.
\item The poles of $c_{334}^2$ are the zeros of higher-weight
  Kac determinants, where the $\cW_4$ algebra develops null
  vectors that obstruct closure.
\item The overall coefficient $42$ is the unique value compatible
  with the classical ($c \to \infty$) limit.
\end{enumerate}
\end{remark}


\subsection{The $E_2$ collapse at higher genus}

\begin{lemma}[$E_2$ collapse at higher genus]\label{lem:e2-higher-genus}
For a Koszul chiral algebra $\cA$ on a genus-$g$ curve, the PBW
spectral sequence
$E_1^{p,q}(g) \Rightarrow H^{p+q}(\Omega_g \barB_g(\cA))$
degenerates at~$E_2$.
\end{lemma}

\begin{proof}
The PBW filtration is a filtration by chiral algebras, and the
Koszulness hypothesis ensures that the associated graded is a
Koszul complex in the classical sense.  The $d_r$ differential
for $r \geq 2$ maps between PBW-graded components separated
by~$r$ steps, and the Koszul concentration (all bar cohomology
is in bar-degree~$0$ on the associated graded) forces these maps
to have zero source or zero target.

On each geometric fiber (a fixed smooth or stable curve), the
collision differential is genus-$0$ type (with $d_0^2 = 0$), and
the genus-$g$ quantum corrections live in higher Leray degrees,
contributing only to $d_r$ for $r \geq 2$.  Since $E_2$ is
already concentrated on the diagonal by Koszul concentration,
these higher differentials vanish.
\end{proof}


\end{document}
